\section*{Capítulo \ref{ch:b3:climate-models} --- Modelos climáticos}

\begin{solution}{exo:b3:climate-models:1}
$A = 14.0 - 1.5 = \num{12.5}$ y $B = 0.14 - 0.05 = \num{0.09}$, en
\unit{kcal.cm^{-2}} por mes. Un mes es \qty{2.63e6}{s}, así que
\qty{1}{kcal.cm^{-2}} por mes es $\num{4.184e7}/\num{2.63e6} =
\qty{15.91}{W/m^2}$: $A = \qty{199}{W/m^2}$, $B =
\qty{1.43}{W.m^{-2}.\celsius^{-1}}$. A \qty{15}{\celsius}, $199 + 21.5 =
\qty{220}{W/m^2}$.
\end{solution}

\begin{solution}{exo:b3:climate-models:2}
$(238.0 - 211.9)/1.816 = \qty{14.4}{\celsius}$; con 0.31, $(234.6 -
211.9)/1.816 = \qty{12.5}{\celsius}$. Cada 0.01 de albedo cuesta
\qty{3.4}{W/m^2} y $3.4/1.816 = \qty{1.9}{\celsius}$.
\end{solution}

\begin{solution}{exo:b3:climate-models:3}
$C = 1025\times3992\times60 = \qty{2.46e8}{J.m^{-2}.K^{-1}}$; dividido por
el \qty{3.156e7}{s} de un año, \qty{7.78}{W.yr.m^{-2}.K^{-1}}. Luego $\tau
= 7.78/1.816 = \num{4.3}~\text{años}$, y $7.78/1.33 = \num{5.8}~\text{años}$.
\end{solution}

\begin{solution}{exo:b3:climate-models:4}
ECS $= 3.93/1.33 = \qty{2.95}{\celsius}$; TCR $\approx 3.93/1.95 =
\qty{2.02}{\celsius}$; razón \num{0.68}.
\end{solution}

\begin{solution}{exo:b3:climate-models:5}
$Q = \qty{335}{W/m^2}$. Bola de nieve: $(335\times0.38 - 211.9)/1.816 =
\qty{-46.6}{\celsius}$; cálida: $(335\times0.70 - 211.9)/1.816 =
\qty{12.4}{\celsius}$, todavía por encima de \qty{10}{\celsius}; rampa: $335(0.54 +
0.016\,T) = 211.9 + 1.816\,T$ da $T = \qty{8.7}{\celsius}$. $N' =
\qty{-1.816}{W.m^{-2}.\celsius^{-1}}$ en los dos primeros, estable, y $-1.816
+ 335\times0.016 = \num{+3.54}$ en la rampa, inestable.
\end{solution}

\begin{solution}{exo:b3:climate-models:6}
La rama cálida termina donde su temperatura alcanza \qty{0}{\celsius}:
$0.70\,Q = A$, $Q = \qty{302.7}{W/m^2}$, con una constante solar de
\qty{1211}{W/m^2}, un \qty{11.0}{\percent} por debajo de la actual. La bola de nieve se funde
donde $0.38\,Q = A - 20B$, $Q = \qty{462.1}{W/m^2}$, $S_0 =
\qty{1848}{W/m^2}$. El umbral es $B(T^* - T_F)/(Q_0(1 - \alpha_F))$,
la diferencia entre la temperatura actual $T^*$ y la temperatura a la que el
hielo empieza a extenderse, convertida en radiación solar: $1.816\times14.37/238 =
\qty{11.0}{\percent}$ aquí, $1.816\times4.37/238 = \qty{3.3}{\percent}$ en
el texto.
\end{solution}

\begin{solution}{exo:b3:climate-models:7}
(a)~\qty{-45.5}{\celsius}. (b)~$340\times0.38 + F = 211.9 - 18.2$: $F =
\qty{64.5}{W/m^2}$, 16 veces el forzamiento de una duplicación. (c)~Mientras el
forzamiento dura, el planeta salta a la rama cálida, $(238.0 + 64.5 -
211.9)/1.816 = \qty{50}{\celsius}$; una vez retirado, el planeta se enfría
por la rama cálida hasta \qty{14.4}{\celsius}, que existe con el Sol actual:
no vuelve a la bola de nieve.
\end{solution}

\begin{solution}{exo:b3:climate-models:8}
$\Sigma = 1.80/7.0 + 0.70/90 = \qty{0.2649}{yr^{-1}}$, $\Pi = 1.10\times0.70/
(7.0\times90) = \qty{1.222e-3}{yr^{-2}}$. Las raíces dan $\tau_f =
\num{3.84}~\text{años}$ y $\tau_s = \num{213}~\text{años}$; las aproximaciones,
$7.0/1.80 = \num{3.89}\text{ años}$ y $90\times1.80/(1.10\times0.70) =
\num{210}~\text{años}$. $a_f = 0.597$ exactamente, $1.10/1.80 = 0.611$
aproximadamente. ECS $= 3.93/1.10 = \qty{3.57}{\celsius}$; TCR $\approx
3.93/1.80 = \qty{2.18}{\celsius}$.
\end{solution}

\begin{solution}{exo:b3:climate-models:9}
Denominador $B + \beta = \num{5.555}$. Ecuador: $(340\times1.221\times0.70 -
211.9 + 3.74\times12.81)/5.555 = \qty{22.8}{\celsius}$; polo:
$(340\times0.506\times0.38 - 211.9 + 47.9)/5.555 = \qty{-17.7}{\celsius}$.
El transporte elimina $3.74\times(22.8 - 12.8) = \qty{37}{W/m^2}$ en el
ecuador y lleva $3.74\times30.6 = \qty{114}{W/m^2}$ hasta el polo.
\end{solution}

\begin{solution}{exo:b3:climate-models:10}
(a)~$(1 - x^2)P_2' = 3x - 3x^3$, cuya derivada es $3 - 9x^2 = -6P_2$.
(b)~El término de transporte es $-\mathcal D T_2(-6P_2) = 6\mathcal DT_2P_2$.
El balance $Q(1 + s_2P_2)(1 - \alpha) - A - B(T_0 + T_2P_2) - 6\mathcal
DT_2P_2 = 0$ se cumple para todo $x$ si su parte constante y su parte en $P_2$
se anulan, lo que da $T_0$ y $T_2$. (c)~$T(0) - T(1) = T_2(-\tfrac12 - 1)
= -1.5\,T_2$. Sin transporte $T_2 = 340\times(-0.476)\times0.70/1.816 =
\qty{-62.4}{\celsius}$: una diferencia de \qty{94}{\celsius}. (d)~$T_2 =
-55.3/1.5 = \qty{-36.9}{\celsius}$, así que $B + 6\mathcal D = 113.3/36.9 =
3.07$ y $\mathcal D = \qty{0.21}{W.m^{-2}.\celsius^{-1}}$.
\end{solution}

\begin{solution}{exo:b3:climate-models:11}
(a)~$r = \tfrac12\Sigma\big(1 \pm \sqrt{1 - 4\Pi/\Sigma^2}\big)$; con
$\sqrt{1 - \epsilon} \approx 1 - \epsilon/2$, las raíces son $\Sigma -
\Pi/\Sigma \approx \Sigma$ y $\Pi/\Sigma$. (b)~Con $T_d = 0$, $C\,\mathrm
dT/\mathrm dt = kt - (\gamma - \lambda)T$, cuya solución desde el reposo es $T =
k(t - \tau_f)/(\gamma - \lambda) + k\tau_f\mathrm e^{-t/\tau_f}/(\gamma -
\lambda)$, $\tau_f = C/(\gamma - \lambda)$. Para $t \gg \tau_f$, $T \approx
kt/(\gamma - \lambda)$, y en la duplicación, $kt = F_{2\times}$: TCR
$\approx F_{2\times}/(\gamma - \lambda)$. (c)~$\Sigma = 1.95/8.1 +
0.62/110 = \qty{0.2464}{yr^{-1}}$, $\Pi = 1.33\times0.62/(8.1\times110) =
\qty{9.25e-4}{yr^{-2}}$; $\tau_f = \num{4.12}~\text{años}$ y $\tau_s =
\num{262}~\text{años}$ (aproximadamente $1/\Sigma = \num{4.06}~\text{años}$ y
$\Sigma/\Pi = \num{266}\text{ años}$); $4\Pi/\Sigma^2 = 0.061$.
\end{solution}

\begin{solution}{exo:b3:climate-models:12}
(a)~Media \qty{3.50}{\celsius}, desviación típica \qty{0.81}{\celsius}.
(b)~ECS $= 1.88 + 1.86\,X$; correlación \num{0.992}; desviación
típica residual \qty{0.11}{\celsius}. (c)~$1.88 + 1.86\times0.7 =
\qty{3.18}{\celsius}$, con una dispersión $\sqrt{(1.86\times0.1)^2 + 0.11^2} =
\qty{0.22}{\celsius}$, alrededor de una cuarta parte de la dispersión del conjunto.
(d)~Ningún mecanismo físico vincula $X$ con las retroalimentaciones; una $X$ hallada
probando muchos observables, de modo que la correlación pueda ser fruto del azar; un sesgo
compartido por los ocho modelos; modelos de la misma familia contados como
independientes.
\end{solution}

\begin{solution}{pb:b3:climate-models:1}
\textbf{1.} $Q_0 = 1360/4 = \qty{340}{W/m^2}$: el planeta intercepta la
luz solar sobre un disco de área $\pi R^2$ y la reparte sobre una esfera de área
$4\pi R^2$.

\textbf{2.} $340\times0.70 = \qty{238.0}{W/m^2}$; $211.9 + 1.816\times15 =
\qty{239.1}{W/m^2}$.

\textbf{3.} $(238.0 - 211.9)/1.816 = \qty{14.4}{\celsius}$.

\textbf{4.} $(129.2 - 211.9)/1.816 = \qty{-45.5}{\celsius}$, por debajo de
\qty{-10}{\celsius}.

\textbf{5.} $340(0.54 + 0.016\,T) = 211.9 + 1.816\,T$: $T =
(183.6 - 211.9)/(1.816 - 5.44) = \qty{7.8}{\celsius}$.

\textbf{6.} $(238.0/\num{5.670e-8})^{1/4} = \qty{254.5}{K}$,
\qty{-18.6}{\celsius}: \qty{33}{\celsius} por debajo del equilibrio cálido, el
efecto invernadero del planeta, incorporado en $A$ y $B$.

\textbf{7.} $N' = -B = \qty{-1.816}{W.m^{-2}.\celsius^{-1}}$ en el
estado de bola de nieve y el equilibrio cálido, ambos estables; $-1.816 + 5.44 =
\num{+3.62}$ en el intermedio, inestable.

\textbf{8.} $\tau = 8.1/3.62 = \num{2.2}\text{ años}$.

\textbf{9.} $8.1/1.816 = \num{4.5}\text{ años}$.

\textbf{10.} $3.93/1.816 = \qty{2.2}{\celsius}$, por debajo de la mejor
estimación evaluada de \qty{3}{\celsius} y por debajo de su rango \emph{probable}, de 2.5 a
\qty{4}{\celsius}: la $B$ del modelo, ajustada a lo largo de las latitudes, no es
el parámetro de retroalimentación climática del clima en su conjunto.

\textbf{11.} $\mathrm e^{-t/\tau} = 0.05$ en $t = 3\tau =
\num{13}~\text{años}$. El clima real también calienta el océano profundo, cuya escala temporal
es de siglos.

\textbf{12.} Para $T_F = \qty{10}{\celsius}$; $Q = (211.9 + 18.2)/0.70 =
\qty{328.7}{W/m^2}$.

\textbf{13.} $S_0 = \qty{1315}{W/m^2}$, \qty{3.3}{\percent} por debajo del de hoy.

\textbf{14.} Hacia el estado de bola de nieve, $(328.7\times0.38 - 211.9)/1.816 =
\qty{-47.9}{\celsius}$.

\textbf{15.} $Q = (211.9 - 18.2)/0.38 = \qty{509.8}{W/m^2}$, $S_0 =
\qty{2039}{W/m^2}$, \qty{50}{\percent} por encima del de hoy.

\textbf{16.} Por la rama cálida hasta \qty{10}{\celsius} con
\qty{1315}{W/m^2}, una caída hasta \qty{-48}{\celsius}, luego de vuelta por la
rama de la bola de nieve hasta \qty{-45.5}{\celsius} con \qty{1360}{W/m^2}: el planeta
permanece congelado. Solo por encima de \qty{2039}{W/m^2} saltaría de nuevo a la
rama cálida.

\textbf{17.} $\bar\alpha = 0.30\times0.9738 + 0.62\times0.0262 =
\num{0.3084}$; $\bar T = (340\times0.6916 - 211.9)/1.816 =
\qty{12.80}{\celsius}$.

\textbf{18.} $T_c = (340\times0.5643\times0.50 - 211.9 + 3.739\times12.80)/
5.555 = \qty{-12.26}{\celsius}$, \qty{3.5}{\celsius} más cálida que la
reanálisis a \ang{72}.

\textbf{19.} $5.555\times(-12.26) + 211.9\times3.059 = \num{580.1}$, con
$\beta/B = \num{2.059}$.

\textbf{20.} Los denominadores son 1.7053, 1.7473, 1.7565, 1.7488 y
1.7362. Los cocientes $Q/Q_0 = \num{1.0005}$, \num{0.9764}, \num{0.9713},
\num{0.9756} y \num{0.9827}.

\textbf{21.} Cerca de \ang{47.5}, con \num{0.971}. Las líneas de hielo situadas hacia el polo respecto del
mínimo son estables, $\mathrm dQ/\mathrm dx_s > 0$; las situadas hacia el ecuador son
inestables.

\textbf{22.} $Q/Q_0 = 1$ también entre \ang{30} y \ang{40}, cerca de
\ang{30} (\ang{30.2} por interpolación): una línea de hielo inestable. En \ang{0},
$Q/Q_0 = 580.1/(0.6106 + 0.7824)/340 = \num{1.225}$: el estado de bola de nieve se mantiene
para cualquier Sol de hasta un \qty{22}{\percent} más brillante que el de hoy, así que es
posible hoy.

\textbf{23.} \qty{2.9}{\percent}, frente a \qty{3.3}{\percent} del
modelo de dimensión cero y \qty{1.6}{\percent} para Budyko; el hielo se sitúa
cerca de \ang{48} en el umbral, frente a su \ang{50}. El modelo latitudinal
se glacia antes que el global, porque la luz solar es más débil donde
se forma el hielo.

\textbf{24.} El planeta del modelo se glacia ante una caída de
\textbf{\qty[text-series-to-math=true]{2.9}{\percent}} de la
constante solar (\qty{2.87}{\percent} según el cálculo).
\end{solution}
